% =====================================================================================================
\section{Result 2: the posterior finds the true distribution}\label{sec:distribution}

\inwords{Suppose the observations are i.i.d., the response follows a sparse Gaussian linear model, the predictors have light
tails, no small group of predictors is almost exactly collinear, and no single predictor carries an unbounded signal.
Then, for every large enough level $M$, with probability at least $1-n^{-c_{\mathrm H}M}$ the posterior puts mass at
most $n^{-c_{\mathrm H}M}$ on conditional laws of the response farther than $\sqrt{M\log n/n}$ from the true one, in the
Hellinger distance. No
lower bound on the coefficients is needed.}

\why{The posterior describes the response correctly even when some true predictors are too weak to be selected, so
prediction does not depend on perfect selection.}

\paragraph{Setting.} In this section only, the observations are $n$ i.i.d.\ copies $(x^{(i)},y^{(i)})$ of a pair $(x,y)$,
with $x\in\R^{\praw}$ the vector of all observed predictors; $\Dn$ is the data set, $\P^*$ its law, and the method runs on
the predictors that are not constant in the sample. The posterior is not defined if fewer than two predictors remain or
the response is constant in the sample; such data sets are counted in the exceptional event of the theorem. A posterior draw $(S,\beta,\sigma^2)$, mapped back to the original
units, defines the conditional law $\Normal(\bar y+(x-\bar x)^\top\beta,\sigma^2)$ of $y$ given $x$; the law of $x$ is not
estimated. Let $f$ be the joint density of $(x,y)$ with $x$ from its true law and $y\mid x$ from this fitted law, and $f^*$
the true joint density; then $d_H(f,f^*):=(1-\int\sqrt{ff^*})^{1/2}$ is the Hellinger distance between the fitted and the true conditional
laws, averaged over a new $x$. $\Pi(\cdot\mid\Dn)$ is the posterior of $(S,\beta,\sigma^2)$, and
$\zeta_j:=(x_j-\E x_j)/\sd(x_j)$. Fix a \emph{level} $M\ge1$; the \emph{window} $q_n:=\min\{\praw,\lceil C_qM\log
n\rceil\}$, $C_q:=2(s(A+2)+4)/\log2$, is the largest number of predictors the conditions constrain jointly.

\begin{localassumption}[Distribution-level setting]\label{asm:dist}
The integer $s\ge0$ and the constants $A\ge1$, $L\ge1$, $c_\psi>0$ and $\bar b>0$ are fixed.
\begin{asmitems}
  \item\label{asm:dist-model} \emph{Gaussian linear model for i.i.d.\ pairs.} $y=\alpha^*+x^\top\bstar+\varepsilon$ with
    $\varepsilon\sim\Normal(0,\sigs^2)$, $\sigs^2>0$, independent of $x$; $\alpha^*$, $\bstar$, $\sigs$ and the law of $x$ may
    depend on $n$.
  \item\label{asm:dist-sparse} \emph{Sparsity and dimension.} At most $s$ coordinates of $\bstar$ are nonzero, and
    $2\le\praw\le n^A$.
  \item\label{asm:dist-subg} \emph{Sub-Gaussian predictors on sparse directions.} Every predictor has positive variance,
    and $\E\exp(\lambda u^\top\zeta)\le\exp\big(L^2\lambda^2\Var(u^\top\zeta)/2\big)$ for every $\lambda\in\R$ and every
    $u\in\R^{\praw}$ with at most $q_n$ nonzero entries.
  \item\label{asm:dist-nec} \emph{No near-exact collinearity.} $\lmin(\Corr(x_T))\ge c_\psi/n$ for every set $T$ of at
    most $q_n$ predictors.
  \item\label{asm:dist-bs} \emph{Bounded signal.} $\abs{\beta^*_j}\sd(x_j)\le\bar b\,\sigs$ for every $j$.
\end{asmitems}
\end{localassumption}

\begin{theorem}[Distribution-level consistency]\label{thm:dist}
There are constants $c_{\mathrm H}>0$, $M_0$ and $n_0$, depending only on $(s,A,L,c_\psi,\bar b)$, such that the following
holds. Fix a level $M\ge M_0$ and assume Assumption~\ref{asm:dist}, with the window $q_n$ of this $M$. Then for every
$n\ge n_0$ with $\sqrt n/\log n\ge M$,
\begin{equation}\label{eq:dist-i}
  \P^*\Big[\,\Pi\Big(d_H(f,f^*)>\sqrt{M\log n/n}\ \Big|\ \Dn\Big)>n^{-c_{\mathrm H} M}\Big]\ \le\ n^{-c_{\mathrm H} M}.
\end{equation}
\end{theorem}

\asmlist{(1)~Assumption~\ref{asm:dist}\ref{asm:dist-model}: Gaussian linear model [\SA{3.1}]; (2)~\ref{asm:dist-sparse}:
sparsity and $\praw\le n^A$ [\SA{3.1}]; (3)~\ref{asm:dist-subg}: light-tailed predictors [\SA{3.2}, law of the rows];
(4)~\ref{asm:dist-nec}: no near-exact collinearity [\SA{3.2}]; (5)~\ref{asm:dist-bs}: bounded signal [\SA{3.2}, signal
size].}

\status{\textbf{Proved and independently checked}. Registry: Appendix~\ref{app:status}, row 2.1.}

\begin{remark}[On the assumptions]\label{rem:dist-asm}
No beta-min condition and no fixed lower bound on eigenvalues are assumed. Jointly Gaussian predictors satisfy Assumption~\ref{asm:dist}\ref{asm:dist-subg} with $L=1$, and every correlation matrix whose $q_n\times q_n$
principal blocks have smallest eigenvalue at least a fixed $\psi>0$ satisfies Assumption~\ref{asm:dist}\ref{asm:dist-nec}
for all large $n$. Condition~\ref{asm:dist-nec} keeps the Occam factor of each extra predictor in force; without it the
radius grows by one factor $\sqrt{\log(\praw\vee n)}$ (Proposition~\ref{prop:dist}). The same statement holds for every
level $M$ up to $\sqrt n/\log n$, with the same constants, It also carries over to the list of models that a completed run returns,
when that list holds posterior mass at least $n^{-c_{\mathrm H}M/2}$, a condition the run does not check
(Corollary~\ref{cor:dist-output}).
\end{remark}
